\section{Datasets considered}
\label{app:datasets}

The HD criterion rewards a ``large number of additional datasets'', but only datasets that a
final model uses go in the zip, and each one costs collection, processing and report words. So
the candidates were checked against four tests: is it Bitcoin, does it carry labels we can use
honestly, is its licence clear, and does it fit a 16 GB machine and a 19-day schedule. The
datasets in use are DS1--DS9 (Section~\ref{sec:datasets}). The table below lists the eighteen
others, checked on 20 October 2026, with the decision for each. It also repeats DS9, the
optional Ransomwhere set, so its reasoning sits beside the rest.

\textbf{Decisions:} \emph{Fallback} means it replaces a chosen dataset if that one fails;
\emph{Could} means it is on the priority list's Could tier; \emph{Not used} means it fails one of
the four tests; \emph{Won't} means it is ruled out for A2 (Section~\ref{sec:moscow}).

{\small
\begin{longtable}{@{}L{3.1cm}L{4.4cm}L{2.4cm}L{4.4cm}@{}}
\toprule
\textbf{Dataset (source)} & \textbf{What it holds} & \textbf{Licence and access} & \textbf{Decision and reason} \\
\midrule
\endhead
\multicolumn{4}{@{}l}{\textbf{Bitcoin, labelled}} \\
Elliptic++ transactions (Elmougy and Liu 2023) & Elliptic's transactions with 183 features and the same labels; also wallet-level (actor) files with 56 features & None stated; Google Drive link from the GitHub README & \emph{Fallback} for DS3 if Elliptic's Kaggle licence blocks it. The actor files are DS8. \\
Elliptic2 (Bellei et al.\ 2024) & 122,000 labelled subgraphs inside a background graph of 49M node clusters and 196M edges & Code Apache 2.0; data licence not stated; download via Elliptic's website & \emph{Not used.} Subgraph labels need graph neural networks (W-1), and the background graph is too large for 16 GB of RAM. \\
BABD-13 (Xiang et al.\ 2024) & 544,462 labelled addresses of 13 behaviour types, 148 features, July 2019 to May 2021 & Host and licence not confirmed & \emph{Not used} unless a confirmed copy appears. It would be a third classification set with little new to say, and its features need a subgraph tool we would have to rebuild. \\
Ransomwhere (Cable 2024) & Crowdsourced ransomware payment addresses with transactions and family; 5.6 MB JSON & CC BY 4.0; Zenodo & \emph{Could}, as DS9. Its addresses run through our own pipeline, which gives a label check in our own feature space. \\
Showing the Receipts (GitHub, \code{cablej}) & 1,013 likely ransomware payment addresses & CC BY 4.0 & \emph{Not used.} Same purpose as Ransomwhere but smaller; kept as a cross-check if DS9 runs. \\
Ransomware seed addresses (Paquet-Clouston, Haslhofer and Dupont 2019) & Seed address CSV, with an expanded set on Zenodo & No licence file & \emph{Not used.} No licence, and it overlaps DS4's label sources. \\
OFAC sanctioned addresses (GitHub, \code{0xB10C}) & Bitcoin addresses on the US sanctions list, regenerated nightly & Repository licence not stated & \emph{Won't.} A ``sanctioned'' label on real addresses is an identity claim, which A1 risk 16 rules out. \\
Bitcoin Ponzi addresses (Bartoletti, Pes and Serusi 2018) & Addresses of Ponzi schemes with features & Host and size not found & \emph{Not used.} No confirmed download. \\
WalletExplorer entity labels & Exchanges, pools and services found by co-spend clustering & No bulk download or licence & \emph{Not used.} No licence; research reports stale names. \\
\multicolumn{4}{@{}l}{\textbf{Bitcoin, whole chain or graph}} \\
ORBITAAL (Coquidé and Cazabet 2024) & Entity-level temporal graph, January 2009 to January 2021 & CC BY 4.0; Zenodo; 156.9 GB & \emph{Won't} (W-5): size. \\
Bitcoin temporal graph (Schnoering and Vazirgiannis 2024) & 252M entity nodes, 786M edges, 34,098 labelled entities & CC BY 4.0; Figshare; PostgreSQL, about 120 GB & \emph{Won't} (W-5): size. Good future work for address profiles. \\
BigQuery \code{crypto_bitcoin} (Google public dataset) & The whole chain as SQL tables & First 1 TB of queries a month free; needs a Google Cloud project & \emph{Fallback} for the DS2 crawl if Esplora blocks it. \\
Blockchair database dumps & Blocks, transactions and outputs as TSV & Licence not confirmed & \emph{Not used.} Licence unclear; Esplora covers the need. \\
SNAP Bitcoin Alpha (Kumar et al.\ 2016) & 3,783 traders, 24,186 trust ratings from $-10$ to $+10$ & Public download from SNAP & \emph{Not used.} Trust ratings between traders, not transactions. \\
\multicolumn{4}{@{}l}{\textbf{Not Bitcoin}} \\
IBM AML (Altman et al.\ 2023) & Synthetic bank transfers with an ``Is Laundering'' label, in six sizes & Kaggle; licence not confirmed & \emph{Not used.} Bank transfers, not Bitcoin. \\
SAML-D (Oztas et al.\ 2023) & About 9.5M synthetic monitored transactions, 12 features, 28 typologies & Kaggle; licence not confirmed & \emph{Not used.} Same reason. \\
SynthAML (Jensen et al.\ 2023) & 20,000 bank alerts (17\% reported) and over 16M transactions & Figshare; data licence not explicit & \emph{Not used.} Same reason. \\
Ethereum fraud sets (Kaggle, XBlock, Forta) & Labelled Ethereum accounts, for example 9,841 accounts with 22\% flagged & Varies & \emph{Won't} (W-4): outside A1's Bitcoin scope. \\
Card-fraud sets (ULB credit card, PaySim, Sparkov) & Card or mobile payments with fraud labels & Varies & \emph{Not used.} Not Bitcoin. \\
\bottomrule
\end{longtable}
}

\textbf{How the report uses this.} One sentence in the data-collection section, ``Eighteen other
public datasets were assessed against four criteria (Appendix~X)'', plus this table condensed in
the report's appendix. That shows selection, which the HD band and data criterion reward, at no
cost to the word limit. Entries that the report names must be cited (the second list in
Appendix~\ref{app:refs}). Details marked ``not confirmed'' are checked again before they appear
in the report.
